\documentclass[10pt]{article}
\usepackage[utf8]{inputenc}
\usepackage[T1]{fontenc}
\usepackage{amsmath}
\usepackage{amsfonts}
\usepackage{amssymb}
\usepackage[version=4]{mhchem}
\usepackage{stmaryrd}
\usepackage{graphicx}
\usepackage[export]{adjustbox}
\graphicspath{ {./images/} }
\usepackage{CJKutf8}

\begin{document}
$\left(x-\lambda^{-1}\right)^{m}$ повторяются одно и то же число раз; 2$)$ каждый элементарный делитель вида $(x \pm 1)^{2 l}$ повторяется четное число раз.

Лит :: [1] М а ль ц е в А. И., Основы линейной алгебры, ОРТОГОНАЈЬНАЯ СЕТЬ - сеть, у к-рой касательные в нек-рой точке к линиям различных семейств ортогональны. Примеры О. с.: асимптотическая сеть па минимальной поверхности, кривизнь линий сеть.

ОРТОГОНАЛЬНАЯ СИСТЕМА - 1) О. с. в е к т ор о в - множество $\left\{x_{\alpha}\right\}$ венулевых векторов евклидова (гильбертова) пространства со скалярным произведением $(\cdot, \cdot)$ такое, что $\left(x_{\alpha}, x_{\beta}\right)=0$ при $\alpha \neq \beta$. Если при этом норма каждого вектора равна единице, то система $\left\{x_{\alpha}\right\}$ наз. о р т о н о р и и р в а н н о й. Полная О. с. $\left\{x_{\alpha}\right\}$ наз. о р тогон алы ны м (ортонор ми-

\begin{center}
\includegraphics[max width=\textwidth]{2023_07_08_6822d4b68da5a18c21bbg-1}
\end{center}

\begin{enumerate}
  \setcounter{enumi}{1}
  \item О.с. ко о р д и н а т-система координат, в к-рой координатные линии (или поверхности) пересекаются под прямым углом. О. с. координат сушествуют в любом евкјидовом пространстве, но, вообще говоря, не существуют в произвольном пространстве. В двумерном гладком афффинном пространстве О. с. всегда можно ввести по крайней мере в достаточно малой окрестности каждой точки. Иногда возможно введеніе О. с. координат в целом. В О. с. метрич. тензор $g_{i j}$ диагонален; диагональные компоненты $g_{i i}$ прьнято наз. коэффицие т т а м Л Л а е. Ламе коэффициент О. с. в пространстве выражағтся формулами
\end{enumerate}

\[
\begin{aligned}
L_{u} & =\sqrt{(\partial x / \partial u)^{2}+(\partial y / \partial u)^{2}+(\partial z / \partial u)^{2}}, \\
L_{v} & =\sqrt{(\partial x / \partial v)^{2}+(\partial y / \partial v)^{2}+(\partial z / \partial v)^{2}}, \\
L_{w} & =\sqrt{(\partial x / \partial w)^{2}+(\partial y / \partial w)^{2}+(\partial z / \partial w)^{2}},
\end{aligned}
\]

где $x, y$ и $z$ - декартовы прямоугольные координаты. Черев коэффициенты Ламе выражаются әлемент длины:

\[
d s=\sqrt{L_{u}^{2} d u^{2}+L_{v}^{2} d v^{2}+L_{w}^{2} d w^{2}}
\]

элемент площади поверхности:

$d \sigma=\sqrt{\left(L_{u} L_{v} d u d v\right)^{2}+\left(L_{u} L_{w} d u d w\right)^{2}+\left(L_{v} L_{w} d v d w\right)^{2}}$, элемент объема:

\[
d V=L_{u} L_{v} L_{w} d u d v d u,
\]

векторные дифференциальные операции:

\[
\begin{aligned}
& \operatorname{grad}_{u} \varphi=\frac{1}{L_{u}} \frac{\partial \varphi}{\partial u}, \operatorname{grad}_{v} \varphi=\frac{1}{L_{v}} \frac{\partial \varphi}{\partial v}, \operatorname{grad}_{w} \varphi=\frac{1}{L_{w}} \frac{\partial \varphi}{\partial w}, \\
& \operatorname{div} \boldsymbol{a}=\frac{1}{L_{u} L_{v} L_{w}}\left[\frac{\partial}{\partial u}\left(a_{u} L_{v} L_{w}\right)+\frac{\partial}{\partial v}\left(a_{v} L_{u} L_{w}\right)+\right. \\
& \left.+\frac{\partial}{\partial w}\left(a_{w} L_{u} L_{v}\right)\right] \\
& \operatorname{rot}_{u} \boldsymbol{a}=\frac{1}{L_{v} L_{w}}\left[\frac{\partial}{\partial v}\left(a_{w} L_{w}\right)-\frac{\partial}{\partial w} i\left(a_{v} L_{v}\right)\right], \\
& \operatorname{rot}_{v} \boldsymbol{a}=\frac{1}{L_{u} L_{w}}\left[\frac{\partial}{\partial w}\left(a_{u} L_{u}\right)-\frac{\partial}{\partial u}\left(a_{w} L_{w}\right)\right] \text {. } \\
& \operatorname{rot}_{w} \boldsymbol{a}=\frac{1}{L_{u} L_{\imath^{\prime}}}\left[\frac{\partial}{\partial u}\left(a_{v} L_{v}\right)-\frac{\partial}{\partial v}\left(a_{u} L_{u}\right)\right], \\
& \Delta \varphi=\frac{1}{L_{u} L_{v} L_{w^{\prime}}}\left[\frac{\partial}{\partial u}\left(\frac{L_{v} L_{w}}{L_{u}} \frac{\partial \varphi}{\partial u}\right)+\frac{\partial}{\partial v}\left(\frac{L_{u} L_{w}}{L_{v}} \frac{\partial \varphi}{\partial v}\right)+\right. \\
& \left.+\frac{\partial}{\partial w}\left(\frac{L_{u} L_{\tilde{v}}}{L_{w}} \frac{\partial \varphi}{\partial w}\right)\right] \text {. }
\end{aligned}
\]

Наиболее часто используемые О. с. координат: на плоскости - декартовы, полярные, эллиптические, параболические; в пространстве - сферические, цилиндрические, параболоидальные, бицилиндрические, биполярные. д. д. Соколов.

\begin{enumerate}
  \setcounter{enumi}{2}
  \item О. с. ф у н к ц и й - конечная или счетная система $\left\{\varphi_{i}(x)\right\}$ функций, принадлежащих пространству $L^{2}(X, S, \mu)$ І удовлетворяющих условиям
\end{enumerate}

\[
\bigcup_{X} \psi_{i}(x) \bar{\varphi}_{j}(x) d \mu(x)=\left\{\begin{array}{lr}
0 & \text { при } i \neq j, \\
\lambda_{i}>0 & \text { при } i=j .
\end{array}\right.
\]

Если $\lambda_{i}=1$ для всех $i$, то система наз. о р т о н о рм и р о в а н в о й. При этом предполагается, что мера $\mu(x)$, определенная на $\sigma$-алгебре $S$ подмножеств множества $X$, счетно аддитивна, полна и имеет счетную базу. Это определение О. с. включает все рассматриваемые в современном анализе 0. с.; они получаются при различных конкретных реализация $\mathrm{x}$ пространства с мерой $(X, S, \mu)$.

Наибольший интерес представляют полные ортонормированные системы $\left\{\varphi_{n}(x)\right\}$, обладающие тем свойством, что для любой функции $f(x) \in L^{2}(X, S, \mu)$ существует единственный ряд $\sum_{c_{n}} \varphi_{n}(x)$, сходящийся к $f(x)$ в метрике простравства $L^{2}(X, S, \mu)$, при этом коэффициенты $c_{n}$ определяются формулами Фурье

\[
c_{n}=\int_{X} \bar{\varphi}_{n} d \mu
\]

Такие системы суцествуют в силу сепарабельности пространства $L^{2}(X, S, \mu)$. Универсальный способ построения полных ортонормированных систем дает метод ортоговализации Шмидта. Для этого достаточно применить его к нек-рой полной в $L^{2}(S, X, \mu)$ системе линейно независимых функций.

В теорий ортогональных рядов в основном рассматриваются О. с. простра\begin{CJK}{UTF8}{mj}甘\end{CJK}ства $L^{2}[a, b]$ (тот частный случай, когда $X=[a, b], S$ - система мвожеств, измеримых по Лебегу, и $\mu$ - мера Лебега). Многие теоремы о сходимости или суммируемости рядов $\sum a_{n} \varphi_{n}(x), \sum a_{n}^{2}<\infty$, по общим О. с. $\left\{\varphi_{n}(x)\right\}$ пространства $L^{2}[a, b]$ верны и для рядов по ортонормированным системам пространства $L^{2}(X, S, \mu)$. Вместе с тем в этом частном случае построены пнтересные конкретные О.с., обладающие теми или ивыми хорошими свойствами. Таковы, например, системы Хаара, Радемахера, Уолша-Пәли, Франклина.

\begin{enumerate}
  \item С и с т е м а $\mathrm{X}$ a a p a $\left\{\chi_{n}(x)\right\}_{n=1}^{\infty}: \quad \chi_{1}(x)=1$, $x \in[0,1]$
\end{enumerate}

\[
\chi_{m}(x)=\left\{\begin{array}{c}
\sqrt{2^{n}} \text { при } x \in\left(\frac{2 k-2}{2^{n+1}}, \frac{2 k-1}{2^{n+1}}\right), \\
-\sqrt{2^{n}} \text { при } x \in\left(\frac{2 k-1}{2^{n+1}}, \frac{2 k}{2^{n+1}}\right), \\
0 \text { в остальных точках отрезка }[0,1],
\end{array}\right.
\]

где $m=2^{n}+k, 1 \leqslant k \leqslant 2^{n}, m=2,3, \ldots$. Ряды по системе Хаара представляют типичный пример мартингалов и для них верны общие теоремы из теорий мартингалов. Кроме того, система $\left\{\chi_{n}(x)\right\}_{n=1}^{\infty}$ является базисом в $L^{p}[0,1], p \geqslant 1$, и ряд Фурье по системе Хаара любой пнтегрируемой функции почти всюду сходится.

\begin{enumerate}
  \setcounter{enumi}{1}
  \item С ил с те м а $\mathrm{P}$ а де м а хе р a $\left\{r_{n}(x)\right\}_{n=1}^{\infty}$ :
\end{enumerate}

\[
r_{n}(x)=\operatorname{sign} \sin 2^{n+1} \pi x, \quad x \in[0,1],
\]

представляет собой важный пример О. с. независимых функций и имеет применения как в теории вероятностей, так и в теории ортогональных и общих функциональных рядов.

\begin{enumerate}
  \setcounter{enumi}{2}
  \item Си с тема уоліп а - П эли $\left\{W_{n}(x)\right\}_{n=1}^{\infty}$ определяется через функции Радемахера:
\end{enumerate}

\[
W_{0}(x)=1, W_{n}(x)=\prod_{k=0}^{m}\left[r_{k}(x)\right]^{q}, x \in[0,1],
\]

где числа $m$ и $q_{k}$ определяются из двоичного разложения чіІла $n$ :

\[
n=\sum_{k=0}^{m} q_{k} 2^{k}
\]


\end{document}